% =====================================================================
%  UNIDAD II --- Lenguajes, Frameworks y Seguridad del Servidor
%  Parte 1: Capitulos 5 y 6 (lenguajes e introduccion + CRUD sin framework)
%  Reutiliza el preambulo de main.tex.
% =====================================================================
\part{Lenguajes, Frameworks y Seguridad del Servidor}

% =====================================================================
\chapter{Introducción a los lenguajes del lado del servidor}
\label{cap:lenguajes}

\epigraph{«Conocer cinco lenguajes del servidor no es saber cinco
sintaxis: es entender que todos resuelven el mismo ciclo
petición--respuesta de maneras distintas.»}{--- \textit{Principio
pedagógico de la asignatura}}

\section*{Resultado de aprendizaje}
\addcontentsline{toc}{section}{Resultado de aprendizaje}
Al terminar este capítulo, el estudiante reconoce el papel del lado del
servidor en una aplicación web, identifica cinco tecnologías de servidor
---Perl, PHP, Python, ASP.NET Core (C\#) y Java para la web--- y escribe
en cada una un programa mínimo «Hola Mundo» que responde por HTTP,
conociendo además la estructura de archivos que genera cada proyecto.

\section{El lado del servidor y el ciclo petición--respuesta}
\label{sec:cap5:ciclo}

Cuando el navegador solicita una página dinámica, un programa en el
servidor recibe la petición HTTP, ejecuta lógica (consulta datos,
calcula, decide) y devuelve una respuesta ---casi siempre HTML5--- que
el navegador presenta. Ese programa puede estar escrito en muchos
lenguajes; este libro recorre cinco de uso real.

\begin{cajadefinicion}{Concepto --- Cliente y servidor}
El \emph{cliente} (navegador) ejecuta HTML5, CSS3 y JavaScript y se
ocupa de la presentación. El \emph{servidor} ejecuta el lenguaje de
turno (Perl, PHP, Python, C\# o Java) y se ocupa de la lógica y de los
datos. La frontera entre ambos es el protocolo HTTP.
\end{cajadefinicion}

\begin{cajaadvertencia}{Punto de partida que asume este libro}
Se asume que el lector ya programa en \textbf{Java} (Java SE: clases,
métodos, colecciones), pero \textbf{no} conoce Perl, PHP, Python ni C\#,
ni el desarrollo web con Java. Por eso cada lenguaje se introduce desde
cero, y Java para la web se presenta al final, como puente desde lo que
ya sabe.
\end{cajaadvertencia}

El IDE de referencia es \textbf{IntelliJ IDEA Ultimate} para Perl, PHP,
Python y Java (mediante sus plugins). Para C\# / ASP.NET Core se usa
\textbf{Visual Studio Code}, porque IntelliJ IDEA no ofrece soporte de
C\#.

\section{Perl}
\label{sec:cap5:perl}

\begin{cajahistoria}{Historia --- Perl}
Creado por Larry Wall en 1987, Perl fue durante los años noventa el
lenguaje dominante de la web a través de CGI. Su lema «hay más de una
manera de hacerlo» y su potencia con expresiones regulares lo hicieron
omnipresente en scripts de servidor. Hoy convive con lenguajes más
nuevos, pero sigue vigente; la versión de referencia es Perl~5.42.
\end{cajahistoria}

Perl genera páginas escribiendo la respuesta HTTP por la salida
estándar. El siguiente «Hola Mundo» es un script CGI: imprime una
cabecera \texttt{Content-Type} y luego el HTML5.

{\small\textbf{Crear en:} \texttt{cgi-bin/hola.pl}}
\begin{lstlisting}[style=perl,caption={hola.pl --- Hola Mundo en Perl (CGI)},label={lst:c5-perl-hola}]
#!/usr/bin/perl
use strict;
use warnings;

print "Content-Type: text/html; charset=UTF-8\n\n";
print "<!DOCTYPE html>\n";
print "<html lang=\"es\"><head><meta charset=\"UTF-8\">";
print "<title>Hola</title></head>";
print "<body><h1>Hola Mundo desde Perl</h1></body></html>";
\end{lstlisting}

\begin{cajaejemplo}{Estructura del proyecto Perl}
\begin{lstlisting}[style=shell,caption={Estructura minima de un proyecto Perl/CGI},label={lst:c5-perl-tree}]
hola-perl/
|-- cgi-bin/
|   `-- hola.pl              (programador)
`-- public/
    `-- estilos.css          (hoja base de la Unidad I)
\end{lstlisting}
\end{cajaejemplo}

Se ejecuta con un servidor que soporte CGI (Apache con
\texttt{mod\_cgi}, o el servidor de pruebas \texttt{plackup}). La sintaxis
clave: las variables escalares empiezan con \texttt{\$}, \texttt{use
strict} obliga a declararlas, y \texttt{print} envía texto al cliente.

\section{PHP}
\label{sec:cap5:php}

\begin{cajahistoria}{Historia --- PHP}
PHP (hoy «PHP: Hypertext Preprocessor») nació en 1994 de la mano de
Rasmus Lerdorf. Diseñado específicamente para la web, se incrusta dentro
del HTML y se ejecuta en el servidor. Es uno de los lenguajes más usados
del mundo (WordPress, Laravel); la versión de referencia es PHP~8.4.
\end{cajahistoria}

En PHP el código va entre \texttt{<?php} y \texttt{?>}, intercalado con
HTML. El intérprete ejecuta el bloque y deja el resto tal cual.

{\small\textbf{Crear en:} \texttt{public/index.php}}
\begin{lstlisting}[style=php,caption={index.php --- Hola Mundo en PHP},label={lst:c5-php-hola}]
<!DOCTYPE html>
<html lang="es">
<head><meta charset="UTF-8"><title>Hola</title>
  <link rel="stylesheet" href="estilos.css"></head>
<body>
  <h1><?php echo "Hola Mundo desde PHP " . phpversion(); ?></h1>
</body>
</html>
\end{lstlisting}

\begin{cajaejemplo}{Estructura del proyecto PHP}
\begin{lstlisting}[style=shell,caption={Estructura minima de un proyecto PHP},label={lst:c5-php-tree}]
hola-php/
`-- public/
    |-- index.php            (programador)
    `-- estilos.css          (hoja base de la Unidad I)
\end{lstlisting}
\end{cajaejemplo}

Se ejecuta con el servidor incorporado: \texttt{php -S localhost:8000 -t
public}. La sintaxis clave: las variables empiezan con \texttt{\$}, las
cadenas se concatenan con el punto (\texttt{.}) y \texttt{echo} envía
texto al cliente.

\section{Python}
\label{sec:cap5:python}

\begin{cajahistoria}{Historia --- Python}
Creado por Guido van Rossum en 1991, Python destaca por su legibilidad.
En la web se usa con frameworks (Django, Flask), pero también se puede
servir HTTP con la biblioteca estándar mediante WSGI ---la interfaz
estándar entre un servidor y una aplicación Python---. La versión de
referencia es Python~3.13.
\end{cajahistoria}

Para mantener el capítulo «sin framework», el «Hola Mundo» usa WSGI puro:
una función que recibe el entorno de la petición y devuelve la respuesta.

{\small\textbf{Crear en:} \texttt{app.py}}
\begin{lstlisting}[style=python,caption={app.py --- Hola Mundo en Python (WSGI)},label={lst:c5-py-hola}]
from wsgiref.simple_server import make_server


def app(environ, start_response):
    cuerpo = ("<!DOCTYPE html><html lang='es'><head>"
              "<meta charset='UTF-8'><title>Hola</title></head>"
              "<body><h1>Hola Mundo desde Python</h1></body></html>")
    datos = cuerpo.encode("utf-8")
    start_response("200 OK", [
        ("Content-Type", "text/html; charset=UTF-8"),
        ("Content-Length", str(len(datos))),
    ])
    return [datos]


if __name__ == "__main__":
    with make_server("", 8000, app) as servidor:
        print("Escuchando en http://localhost:8000")
        servidor.serve_forever()
\end{lstlisting}

\begin{cajaejemplo}{Estructura del proyecto Python}
\begin{lstlisting}[style=shell,caption={Estructura minima de un proyecto Python/WSGI},label={lst:c5-py-tree}]
hola-python/
|-- app.py                   (programador)
`-- static/
    `-- estilos.css          (hoja base de la Unidad I)
\end{lstlisting}
\end{cajaejemplo}

Se ejecuta con \texttt{python app.py}. La sintaxis clave: la indentación
define los bloques, no hay llaves; \texttt{def} declara funciones y las
cadenas se codifican a bytes antes de enviarse.

\section{ASP.NET Core (C\#)}
\label{sec:cap5:asp}

\begin{cajahistoria}{Historia --- ASP.NET Core}
Microsoft lanzó ASP.NET en 2002; en 2016 lo reescribió como
\textbf{ASP.NET Core}, multiplataforma y de código abierto. C\# es su
lenguaje, de tipado fuerte y sintaxis cercana a Java. La versión de
referencia es \textbf{.NET~10 (LTS)} con \textbf{C\#~14} y ASP.NET
Core~10.
\end{cajahistoria}

ASP.NET Core usa el modelo de \emph{Minimal API}: en pocas líneas se
define una ruta y su respuesta. El proyecto se crea con la CLI de .NET.

{\small\textbf{Crear el proyecto (terminal de VS Code):}}
\begin{lstlisting}[style=shell,caption={Creacion del proyecto ASP.NET Core},label={lst:c5-asp-new}]
dotnet new web -o HolaAsp
cd HolaAsp
dotnet run
\end{lstlisting}

{\small\textbf{Editar:} \texttt{Program.cs}}
\begin{lstlisting}[style=cs,caption={Program.cs --- Hola Mundo en ASP.NET Core},label={lst:c5-asp-hola}]
var builder = WebApplication.CreateBuilder(args);
var app = builder.Build();

app.MapGet("/", () => Results.Content(
    "<!DOCTYPE html><html lang=\"es\"><head>" +
    "<meta charset=\"UTF-8\"><title>Hola</title></head>" +
    "<body><h1>Hola Mundo desde ASP.NET Core</h1></body></html>",
    "text/html"));

app.Run();
\end{lstlisting}

\begin{cajaejemplo}{Estructura del proyecto ASP.NET Core}
\begin{lstlisting}[style=shell,caption={Estructura creada por dotnet new web},label={lst:c5-asp-tree}]
HolaAsp/
|-- Program.cs               (dotnet; lo edita el programador)
|-- HolaAsp.csproj           (dotnet: proyecto y dependencias)
|-- appsettings.json         (dotnet: configuracion)
|-- Properties/
|   `-- launchSettings.json   (dotnet: perfiles de ejecucion)
`-- wwwroot/                  (programador: estaticos, estilos.css)
\end{lstlisting}
\end{cajaejemplo}

La sintaxis clave: C\# es de tipado fuerte; \texttt{var} infiere el tipo;
los \emph{lambdas} (\texttt{() => \ldots}) definen el manejador de la
ruta. Se ejecuta con \texttt{dotnet run}.

\section{Java para la Web}
\label{sec:cap5:javaweb}

Como el lector ya domina Java SE, aquí solo se tiende el puente hacia la
web. En Java, un \textbf{servlet} es una clase que atiende peticiones
HTTP: hereda de \texttt{HttpServlet} y sobrescribe \texttt{doGet}. Se
ejecuta dentro de un servidor de aplicaciones; el de referencia es
\textbf{Apache Tomcat~11} (Jakarta Servlet~6.1), con \textbf{Java~21}.

{\small\textbf{Crear en:} \texttt{src/main/java/ec/edu/uteq/HolaServlet.java}}
\begin{lstlisting}[style=java,caption={HolaServlet.java --- Hola Mundo en Java web},label={lst:c5-java-hola}]
package ec.edu.uteq;

import jakarta.servlet.annotation.WebServlet;
import jakarta.servlet.http.HttpServlet;
import jakarta.servlet.http.HttpServletRequest;
import jakarta.servlet.http.HttpServletResponse;
import java.io.IOException;

@WebServlet("/hola")
public class HolaServlet extends HttpServlet {
    @Override
    protected void doGet(HttpServletRequest req, HttpServletResponse resp)
            throws IOException {
        resp.setContentType("text/html; charset=UTF-8");
        resp.getWriter().println(
            "<!DOCTYPE html><html lang=\"es\"><head>"
            + "<meta charset=\"UTF-8\"><title>Hola</title></head>"
            + "<body><h1>Hola Mundo desde Java (Servlet)</h1></body></html>");
    }
}
\end{lstlisting}

\begin{cajaejemplo}{Estructura del proyecto Java web (Maven --- war)}
\begin{lstlisting}[style=shell,caption={Estructura minima de un proyecto Java web},label={lst:c5-java-tree}]
hola-java/
|-- pom.xml                              (programador: dependencias)
`-- src/main/
    |-- java/ec/edu/uteq/
    |   `-- HolaServlet.java             (programador)
    `-- webapp/
        |-- WEB-INF/web.xml              (IntelliJ)
        `-- estilos.css                  (hoja base de la Unidad I)
\end{lstlisting}
\end{cajaejemplo}

La anotación \texttt{@WebServlet("/hola")} mapea la URL; no hace falta
\texttt{web.xml} para registrar el servlet. El detalle completo de
JSP/Servlets/JSTL/EL llega en el Capítulo~\ref{cap:crud}.

\begin{cajacompara}{Las cinco tecnologías de un vistazo}
\small
\begin{tabular}{>{\bfseries}l l l}
\toprule
Tecnología & Cómo responde HTTP & Se ejecuta con \\
\midrule
Perl   & \texttt{print} a la salida (CGI) & Apache/\texttt{plackup} \\
PHP    & Código incrustado en HTML        & \texttt{php -S} \\
Python & Función WSGI                     & \texttt{python app.py} \\
C\#/ASP & \emph{Minimal API} (\texttt{MapGet}) & \texttt{dotnet run} \\
Java   & Servlet (\texttt{doGet})         & Tomcat 11 \\
\bottomrule
\end{tabular}
\end{cajacompara}

% =====================================================================
\chapter{CRUD sin framework: acceso a datos a mano}
\label{cap:crud}

\epigraph{«Antes de dejar que un framework escriba el SQL por ti,
escríbelo tú: solo así sabrás qué hace ---y qué deshacer--- cuando
falle.»}{--- \textit{Principio pedagógico de la asignatura}}

\section*{Resultado de aprendizaje}
\addcontentsline{toc}{section}{Resultado de aprendizaje}
Al terminar este capítulo, el estudiante implementa una aplicación CRUD
(crear, leer, actualizar y eliminar) sobre la tabla \texttt{estudiantes}
en los cinco lenguajes, escribiendo a mano la conexión y el SQL
parametrizado, sin emplear ningún framework ni ORM. En cada caso conoce
la estructura del proyecto y reutiliza la hoja \texttt{estilos.css} de la
Unidad~I.

\section{El modelo de datos común}
\label{sec:cap6:modelo}

Las cinco implementaciones operan sobre la misma base PostgreSQL
\texttt{appweb2026ppa} y la misma tabla. Se crea una sola vez.

{\small\textbf{Crear en:} \texttt{db/schema.sql}}
\begin{lstlisting}[style=sql,caption={Esquema compartido de la tabla estudiantes},label={lst:c6-schema}]
CREATE TABLE IF NOT EXISTS estudiantes (
    id      SERIAL       PRIMARY KEY,
    nombre  VARCHAR(100) NOT NULL,
    correo  VARCHAR(120) NOT NULL,
    carrera VARCHAR(80)  NOT NULL,
    creado  TIMESTAMP    NOT NULL DEFAULT NOW()
);
\end{lstlisting}

\begin{cajaadvertencia}{Credenciales y SQL parametrizado}
Servidor \texttt{localhost:5432}, base \texttt{appweb2026ppa}, usuario
\texttt{postgres}, contraseña \texttt{P@\$7Gr3\$QL}. En las cinco
tecnologías los datos del usuario viajan como \textbf{parámetros}
enlazados, nunca concatenados en la cadena SQL: esa es la defensa contra
la inyección SQL.
\end{cajaadvertencia}

\section{Nuevo CSS para el listado}
\label{sec:cap6:css}

La hoja \texttt{estilos.css} de la Unidad~I ya define el reinicio, la
tipografía y la estructura. Para las tablas del CRUD solo se agregan
estas reglas \emph{nuevas}; el resto se hereda.

{\small\textbf{Agregar a:} \texttt{estilos.css}}
\begin{lstlisting}[style=css,caption={Reglas nuevas para la tabla del CRUD},label={lst:c6-css}]
.tabla-crud { width: 100%; border-collapse: collapse; margin: 1rem 0; }
.tabla-crud th, .tabla-crud td {
  border: 1px solid var(--borde); padding: .5rem .75rem; text-align: left;
}
.tabla-crud th { background: var(--primario); color: #fff; }
.tabla-crud tr:nth-child(even) { background: var(--fondo-alt); }
.acciones a { margin-right: .5rem; }
\end{lstlisting}

% ---------------------------------------------------------------------
\section{CRUD en Perl con DBI}
\label{sec:cap6:perl}

Perl accede a PostgreSQL con el módulo \textbf{DBI} y el driver
\textbf{DBD::Pg}. La conexión se centraliza en una subrutina; cada
operación usa \texttt{prepare} y \texttt{execute} con marcadores
\texttt{?}.

{\small\textbf{Crear en:} \texttt{lib/EstudianteDao.pm} --- parte 1 de 2 (conexion y lectura)}
\begin{lstlisting}[style=perl,caption={EstudianteDao.pm --- parte 1: conexion y listar},label={lst:c6-perl-a}]
package EstudianteDao;
use strict;
use warnings;
use DBI;

sub conectar {
    my $dsn = "dbi:Pg:dbname=appweb2026ppa;host=localhost;port=5432";
    return DBI->connect($dsn, "postgres", 'P@$7Gr3$QL',
        { RaiseError => 1, AutoCommit => 1 });
}

sub listar {
    my $dbh = conectar();
    my $sth = $dbh->prepare(
        "SELECT id, nombre, correo, carrera FROM estudiantes ORDER BY id");
    $sth->execute();
    my @lista;
    while (my $fila = $sth->fetchrow_hashref) { push @lista, $fila; }
    $dbh->disconnect;
    return \@lista;
}
\end{lstlisting}

{\small\textbf{Continúa en:} \texttt{lib/EstudianteDao.pm} --- parte 2 de 2 (escritura)}
\begin{lstlisting}[style=perl,caption={EstudianteDao.pm --- parte 2: insertar, actualizar, eliminar},label={lst:c6-perl-b}]
sub insertar {
    my ($nombre, $correo, $carrera) = @_;
    my $dbh = conectar();
    $dbh->do("INSERT INTO estudiantes (nombre, correo, carrera) "
           . "VALUES (?, ?, ?)", undef, $nombre, $correo, $carrera);
    $dbh->disconnect;
}

sub actualizar {
    my ($id, $nombre, $correo, $carrera) = @_;
    my $dbh = conectar();
    $dbh->do("UPDATE estudiantes SET nombre=?, correo=?, carrera=? "
           . "WHERE id=?", undef, $nombre, $correo, $carrera, $id);
    $dbh->disconnect;
}

sub eliminar {
    my ($id) = @_;
    my $dbh = conectar();
    $dbh->do("DELETE FROM estudiantes WHERE id=?", undef, $id);
    $dbh->disconnect;
}
1;
\end{lstlisting}

El script CGI usa el DAO y genera el listado HTML5 enlazando
\texttt{estilos.css}.

{\small\textbf{Crear en:} \texttt{cgi-bin/lista.pl}}
\begin{lstlisting}[style=perl,caption={lista.pl --- vista de listado},label={lst:c6-perl-vista}]
#!/usr/bin/perl
use strict; use warnings; use lib "../lib"; use EstudianteDao;

print "Content-Type: text/html; charset=UTF-8\n\n";
print "<!DOCTYPE html><html lang=\"es\"><head><meta charset=\"UTF-8\">";
print "<link rel=\"stylesheet\" href=\"/estilos.css\"></head><body>";
print "<h1>Estudiantes</h1><table class=\"tabla-crud\">";
print "<tr><th>ID</th><th>Nombre</th><th>Correo</th><th>Carrera</th></tr>";
for my $e (@{ EstudianteDao::listar() }) {
    print "<tr><td>$e->{id}</td><td>$e->{nombre}</td>";
    print "<td>$e->{correo}</td><td>$e->{carrera}</td></tr>";
}
print "</table></body></html>";
\end{lstlisting}

\begin{cajaejemplo}{Estructura del proyecto CRUD en Perl}
\begin{lstlisting}[style=shell,caption={Estructura del CRUD Perl},label={lst:c6-perl-tree}]
crud-perl/
|-- db/schema.sql                (programador)
|-- lib/
|   `-- EstudianteDao.pm          (programador: acceso a datos)
|-- cgi-bin/
|   |-- lista.pl                  (programador: vista)
|   |-- guardar.pl                (programador: alta/edicion)
|   `-- eliminar.pl               (programador)
`-- public/estilos.css            (hoja base de la Unidad I)
\end{lstlisting}
\end{cajaejemplo}

% ---------------------------------------------------------------------
\section{CRUD en PHP con PDO}
\label{sec:cap6:php}

PHP usa \textbf{PDO} (PHP Data Objects), una capa uniforme de acceso a
datos. La conexión se crea una vez y se reutiliza; las consultas usan
marcadores con nombre (\texttt{:nombre}).

{\small\textbf{Crear en:} \texttt{src/conexion.php}}
\begin{lstlisting}[style=php,caption={conexion.php --- conexion PDO},label={lst:c6-php-cn}]
<?php
function conectar(): PDO {
    $dsn = 'pgsql:host=localhost;port=5432;dbname=appweb2026ppa';
    return new PDO($dsn, 'postgres', 'P@$7Gr3$QL', [
        PDO::ATTR_ERRMODE => PDO::ERRMODE_EXCEPTION,
        PDO::ATTR_DEFAULT_FETCH_MODE => PDO::FETCH_ASSOC,
    ]);
}
\end{lstlisting}

{\small\textbf{Crear en:} \texttt{src/EstudianteDao.php}}
\begin{lstlisting}[style=php,caption={EstudianteDao.php --- CRUD con PDO},label={lst:c6-php-dao}]
<?php
require_once 'conexion.php';

class EstudianteDao {
    public function listar(): array {
        $sql = 'SELECT id, nombre, correo, carrera FROM estudiantes ORDER BY id';
        return conectar()->query($sql)->fetchAll();
    }
    public function insertar(string $nombre, string $correo, string $carrera): void {
        $sql = 'INSERT INTO estudiantes (nombre, correo, carrera) '
             . 'VALUES (:nombre, :correo, :carrera)';
        $stmt = conectar()->prepare($sql);
        $stmt->execute([':nombre'=>$nombre, ':correo'=>$correo, ':carrera'=>$carrera]);
    }
    public function actualizar(int $id, string $nombre, string $correo, string $carrera): void {
        $sql = 'UPDATE estudiantes SET nombre=:nombre, correo=:correo, '
             . 'carrera=:carrera WHERE id=:id';
        $stmt = conectar()->prepare($sql);
        $stmt->execute([':nombre'=>$nombre, ':correo'=>$correo,
                        ':carrera'=>$carrera, ':id'=>$id]);
    }
    public function eliminar(int $id): void {
        $stmt = conectar()->prepare('DELETE FROM estudiantes WHERE id=:id');
        $stmt->execute([':id'=>$id]);
    }
}
\end{lstlisting}

{\small\textbf{Crear en:} \texttt{public/index.php}}
\begin{lstlisting}[style=php,caption={index.php --- vista de listado},label={lst:c6-php-vista}]
<?php require_once '../src/EstudianteDao.php';
$lista = (new EstudianteDao())->listar(); ?>
<!DOCTYPE html>
<html lang="es"><head><meta charset="UTF-8">
  <link rel="stylesheet" href="estilos.css"></head>
<body>
  <h1>Estudiantes</h1>
  <table class="tabla-crud">
    <tr><th>ID</th><th>Nombre</th><th>Correo</th><th>Carrera</th></tr>
    <?php foreach ($lista as $e): ?>
      <tr><td><?= htmlspecialchars($e['id']) ?></td>
          <td><?= htmlspecialchars($e['nombre']) ?></td>
          <td><?= htmlspecialchars($e['correo']) ?></td>
          <td><?= htmlspecialchars($e['carrera']) ?></td></tr>
    <?php endforeach; ?>
  </table>
</body></html>
\end{lstlisting}

\begin{cajaejemplo}{Estructura del proyecto CRUD en PHP}
\begin{lstlisting}[style=shell,caption={Estructura del CRUD PHP},label={lst:c6-php-tree}]
crud-php/
|-- db/schema.sql                (programador)
|-- src/
|   |-- conexion.php             (programador: conexion)
|   `-- EstudianteDao.php        (programador: acceso a datos)
`-- public/
    |-- index.php                (programador: vista)
    `-- estilos.css              (hoja base de la Unidad I)
\end{lstlisting}
\end{cajaejemplo}

La función \texttt{htmlspecialchars} escapa la salida: defensa anti-XSS,
equivalente a \texttt{c:out} en JSTL.

% ---------------------------------------------------------------------
\section{CRUD en Python con psycopg (sin framework)}
\label{sec:cap6:python}

Para mantener el capítulo sin framework, Python accede a PostgreSQL con
el adaptador \textbf{psycopg} (versión 3) y sirve HTTP por WSGI. El DAO
ejecuta SQL parametrizado con marcadores \texttt{\%s}.

{\small\textbf{Instalar:} \texttt{pip install "psycopg[binary]"}}

{\small\textbf{Crear en:} \texttt{dao.py}}
\begin{lstlisting}[style=python,caption={dao.py --- CRUD con psycopg},label={lst:c6-py-dao}]
import psycopg

CADENA = ("host=localhost port=5432 dbname=appweb2026ppa "
          "user=postgres password=P@$7Gr3$QL")


def listar():
    with psycopg.connect(CADENA) as cn, cn.cursor() as cur:
        cur.execute("SELECT id, nombre, correo, carrera "
                    "FROM estudiantes ORDER BY id")
        return cur.fetchall()


def insertar(nombre, correo, carrera):
    with psycopg.connect(CADENA) as cn, cn.cursor() as cur:
        cur.execute("INSERT INTO estudiantes (nombre, correo, carrera) "
                    "VALUES (%s, %s, %s)", (nombre, correo, carrera))


def actualizar(id_est, nombre, correo, carrera):
    with psycopg.connect(CADENA) as cn, cn.cursor() as cur:
        cur.execute("UPDATE estudiantes SET nombre=%s, correo=%s, "
                    "carrera=%s WHERE id=%s",
                    (nombre, correo, carrera, id_est))


def eliminar(id_est):
    with psycopg.connect(CADENA) as cn, cn.cursor() as cur:
        cur.execute("DELETE FROM estudiantes WHERE id=%s", (id_est,))
\end{lstlisting}

La aplicación WSGI construye el listado HTML5 y enlaza
\texttt{estilos.css}.

{\small\textbf{Crear en:} \texttt{app.py}}
\begin{lstlisting}[style=python,caption={app.py --- vista de listado (WSGI)},label={lst:c6-py-app}]
from wsgiref.simple_server import make_server
from html import escape
import dao


def app(environ, start_response):
    filas = "".join(
        f"<tr><td>{e[0]}</td><td>{escape(e[1])}</td>"
        f"<td>{escape(e[2])}</td><td>{escape(e[3])}</td></tr>"
        for e in dao.listar())
    html = (f"<!DOCTYPE html><html lang='es'><head>"
            f"<meta charset='UTF-8'>"
            f"<link rel='stylesheet' href='/estilos.css'></head><body>"
            f"<h1>Estudiantes</h1><table class='tabla-crud'>"
            f"<tr><th>ID</th><th>Nombre</th><th>Correo</th>"
            f"<th>Carrera</th></tr>{filas}</table></body></html>")
    datos = html.encode("utf-8")
    start_response("200 OK", [("Content-Type", "text/html; charset=UTF-8")])
    return [datos]


if __name__ == "__main__":
    with make_server("", 8000, app) as s:
        s.serve_forever()
\end{lstlisting}

\begin{cajaejemplo}{Estructura del proyecto CRUD en Python}
\begin{lstlisting}[style=shell,caption={Estructura del CRUD Python},label={lst:c6-py-tree}]
crud-python/
|-- db/schema.sql                (programador)
|-- dao.py                       (programador: acceso a datos)
|-- app.py                       (programador: vista WSGI)
`-- static/estilos.css           (hoja base de la Unidad I)
\end{lstlisting}
\end{cajaejemplo}

\begin{cajaadvertencia}{Sobre «sin framework» en Python}
Hacer un CRUD web sin ningún framework es inusual: por eso se usa la
biblioteca estándar (WSGI) y \texttt{psycopg} a bajo nivel. En el
Capítulo~\ref{cap:frameworks} se reescribe este mismo CRUD con
\textbf{Django}, que aporta enrutamiento, plantillas y mucho más.
\end{cajaadvertencia}

% ---------------------------------------------------------------------
\section{CRUD en ASP.NET Core con ADO.NET (sin EF)}
\label{sec:cap6:asp}

C\# accede a PostgreSQL con el proveedor \textbf{Npgsql} (ADO.NET). Para
que sea parejo con los demás, no se usa Entity Framework: el SQL se
escribe a mano con parámetros (\texttt{@nombre}).

{\small\textbf{Instalar:} \texttt{dotnet add package Npgsql} (versión 10)}

{\small\textbf{Crear en:} \texttt{Datos/EstudianteDao.cs} --- parte 1 de 2 (conexion y lectura)}
\begin{lstlisting}[style=cs,caption={EstudianteDao.cs --- parte 1: conexion y listar},label={lst:c6-asp-a}]
using Npgsql;

namespace CrudAsp.Datos;

public class EstudianteDao
{
    private const string Cadena =
        "Host=localhost;Port=5432;Database=appweb2026ppa;" +
        "Username=postgres;Password=P@$7Gr3$QL";

    public List<(int, string, string, string)> Listar()
    {
        var lista = new List<(int, string, string, string)>();
        using var cn = new NpgsqlConnection(Cadena);
        cn.Open();
        using var cmd = new NpgsqlCommand(
            "SELECT id, nombre, correo, carrera FROM estudiantes ORDER BY id", cn);
        using var rd = cmd.ExecuteReader();
        while (rd.Read())
            lista.Add((rd.GetInt32(0), rd.GetString(1),
                       rd.GetString(2), rd.GetString(3)));
        return lista;
    }
}
\end{lstlisting}

{\small\textbf{Continúa en:} \texttt{Datos/EstudianteDao.cs} --- parte 2 de 2 (escritura)}
\begin{lstlisting}[style=cs,caption={EstudianteDao.cs --- parte 2: insertar, actualizar, eliminar},label={lst:c6-asp-b}]
    public void Insertar(string nombre, string correo, string carrera)
    {
        using var cn = new NpgsqlConnection(Cadena);
        cn.Open();
        using var cmd = new NpgsqlCommand(
            "INSERT INTO estudiantes (nombre, correo, carrera) " +
            "VALUES (@n, @c, @ca)", cn);
        cmd.Parameters.AddWithValue("n", nombre);
        cmd.Parameters.AddWithValue("c", correo);
        cmd.Parameters.AddWithValue("ca", carrera);
        cmd.ExecuteNonQuery();
    }

    public void Eliminar(int id)
    {
        using var cn = new NpgsqlConnection(Cadena);
        cn.Open();
        using var cmd = new NpgsqlCommand(
            "DELETE FROM estudiantes WHERE id=@id", cn);
        cmd.Parameters.AddWithValue("id", id);
        cmd.ExecuteNonQuery();
    }
}
\end{lstlisting}

El \emph{Minimal API} expone la ruta que arma el listado HTML5.

{\small\textbf{Editar:} \texttt{Program.cs}}
\begin{lstlisting}[style=cs,caption={Program.cs --- vista de listado},label={lst:c6-asp-prog}]
using CrudAsp.Datos;

var app = WebApplication.CreateBuilder(args).Build();
var dao = new EstudianteDao();

app.MapGet("/estudiantes", () =>
{
    var filas = string.Join("", dao.Listar().Select(e =>
        $"<tr><td>{e.Item1}</td><td>{e.Item2}</td>" +
        $"<td>{e.Item3}</td><td>{e.Item4}</td></tr>"));
    var html = "<!DOCTYPE html><html lang=\"es\"><head>" +
        "<meta charset=\"UTF-8\">" +
        "<link rel=\"stylesheet\" href=\"/estilos.css\"></head><body>" +
        "<h1>Estudiantes</h1><table class=\"tabla-crud\">" +
        "<tr><th>ID</th><th>Nombre</th><th>Correo</th><th>Carrera</th></tr>" +
        filas + "</table></body></html>";
    return Results.Content(html, "text/html");
});

app.UseStaticFiles();
app.Run();
\end{lstlisting}

\begin{cajaejemplo}{Estructura del proyecto CRUD en ASP.NET Core}
\begin{lstlisting}[style=shell,caption={Estructura del CRUD ASP.NET Core},label={lst:c6-asp-tree}]
CrudAsp/
|-- Program.cs                   (programador: rutas/vista)
|-- CrudAsp.csproj               (dotnet: proyecto y paquetes)
|-- Datos/
|   `-- EstudianteDao.cs          (programador: acceso a datos)
|-- db/schema.sql                 (programador)
`-- wwwroot/estilos.css           (hoja base de la Unidad I)
\end{lstlisting}
\end{cajaejemplo}

% ---------------------------------------------------------------------
\section{CRUD en Java con JSP, Servlets, JSTL y EL (JDBC)}
\label{sec:cap6:java}

Java reúne aquí las cuatro piezas: \textbf{Servlet} (control),
\textbf{JDBC} (datos), \textbf{JSP} con \textbf{JSTL} y \textbf{EL}
(vista). Es la versión ampliada del CRUD que se introdujo como práctica
integradora; el detalle de cada etiqueta JSTL se trató en su material
específico.

{\small\textbf{Crear en:} \texttt{.../dao/EstudianteDAO.java} --- parte 1 de 2 (conexion y lectura)}
\begin{lstlisting}[style=java,caption={EstudianteDAO.java --- parte 1: conexion y listar},label={lst:c6-java-a}]
package ec.edu.uteq.appweb.dao;

import ec.edu.uteq.appweb.modelo.Estudiante;
import java.sql.*;
import java.util.ArrayList;
import java.util.List;

public class EstudianteDAO {
    private static final String URL =
        "jdbc:postgresql://localhost:5432/appweb2026ppa";

    private Connection cn() throws SQLException {
        return DriverManager.getConnection(URL, "postgres", "P@$7Gr3$QL");
    }

    public List<Estudiante> listar() throws SQLException {
        String sql = "SELECT id, nombre, correo, carrera "
                   + "FROM estudiantes ORDER BY id";
        List<Estudiante> lista = new ArrayList<>();
        try (Connection c = cn();
             PreparedStatement ps = c.prepareStatement(sql);
             ResultSet rs = ps.executeQuery()) {
            while (rs.next())
                lista.add(new Estudiante(rs.getInt("id"),
                    rs.getString("nombre"), rs.getString("correo"),
                    rs.getString("carrera")));
        }
        return lista;
    }
}
\end{lstlisting}

{\small\textbf{Continúa en:} \texttt{.../dao/EstudianteDAO.java} --- parte 2 de 2 (escritura)}
\begin{lstlisting}[style=java,caption={EstudianteDAO.java --- parte 2: insertar, actualizar, eliminar},label={lst:c6-java-b}]
    public void insertar(Estudiante e) throws SQLException {
        String sql = "INSERT INTO estudiantes (nombre, correo, carrera) "
                   + "VALUES (?, ?, ?)";
        try (Connection c = cn();
             PreparedStatement ps = c.prepareStatement(sql)) {
            ps.setString(1, e.getNombre());
            ps.setString(2, e.getCorreo());
            ps.setString(3, e.getCarrera());
            ps.executeUpdate();
        }
    }

    public void eliminar(int id) throws SQLException {
        try (Connection c = cn();
             PreparedStatement ps =
                 c.prepareStatement("DELETE FROM estudiantes WHERE id=?")) {
            ps.setInt(1, id);
            ps.executeUpdate();
        }
    }
}
\end{lstlisting}

La vista JSP recorre la lista con \texttt{c:forEach} y escapa la salida
con \texttt{c:out}.

{\small\textbf{Crear en:} \texttt{WEB-INF/vistas/lista.jsp}}
\begin{lstlisting}[style=html,caption={lista.jsp --- vista con JSTL y EL},label={lst:c6-java-jsp}]
<%@ page contentType="text/html; charset=UTF-8" %>
<%@ taglib prefix="c" uri="jakarta.tags.core" %>
<!DOCTYPE html>
<html lang="es"><head><meta charset="UTF-8">
  <link rel="stylesheet" href="estilos.css"></head>
<body>
  <h1>Estudiantes</h1>
  <table class="tabla-crud">
    <tr><th>ID</th><th>Nombre</th><th>Correo</th><th>Carrera</th></tr>
    <c:forEach var="e" items="${lista}">
      <tr><td>${e.id}</td><td><c:out value="${e.nombre}"/></td>
          <td><c:out value="${e.correo}"/></td>
          <td><c:out value="${e.carrera}"/></td></tr>
    </c:forEach>
  </table>
</body></html>
\end{lstlisting}

\begin{cajaejemplo}{Estructura del proyecto CRUD en Java}
\begin{lstlisting}[style=shell,caption={Estructura del CRUD Java web},label={lst:c6-java-tree}]
crud-java/
|-- pom.xml                                   (programador)
|-- db/schema.sql                             (programador)
`-- src/main/
    |-- java/ec/edu/uteq/appweb/
    |   |-- modelo/Estudiante.java            (programador: dominio)
    |   |-- dao/EstudianteDAO.java            (programador: datos)
    |   `-- controlador/EstudianteServlet.java (programador: control)
    `-- webapp/
        |-- WEB-INF/vistas/lista.jsp          (programador: vista)
        `-- estilos.css                       (hoja base de la Unidad I)
\end{lstlisting}
\end{cajaejemplo}

\section{Síntesis del capítulo}
\label{sec:cap6:sintesis}

Los cinco lenguajes resolvieron el mismo CRUD con la misma idea:
\emph{una conexión propia y SQL parametrizado}. Cambian la sintaxis y el
marcador de parámetro ---\texttt{?} en Perl y Java, \texttt{:nombre} en
PHP, \texttt{\%s} en Python, \texttt{@nombre} en C\#---, pero el patrón
es idéntico. En el Capítulo~\ref{cap:frameworks} se reescriben estos
CRUD con el framework de cada tecnología, donde el ORM generará el SQL
por nosotros.
